\section{One dominant direction: a $q=1$ interval}\label{sec:q1}

When a single direction of the quadratic part of $\hat\theta$ is not negligible relative to its total
uncertainty, as can happen with two groups of firms linked by few movers, condition (b) of
Proposition~\ref{prop:conc} can fail, $\hat\theta$ need not be close to normal, and the interval
$\hat\theta\pm1.96\,\hat\V[\hat\theta]^{1/2}$ should then not be used. For this case KSS (their Sections~6 and~7,
with $q=1$) treat the dominant direction separately. Their interval needs a $2\times2$ covariance matrix, and
the variance estimates of Section~\ref{sec:estimator} can supply it. This section gives the three entries, shows
that the matrix built on them with KSS's own first entry is consistent under conditions parallel to those of
Section~\ref{sec:results} (Proposition~\ref{prop:q1}), shows that the rule used here for the first entry keeps
the coverage that follows from it (Corollary~\ref{cor:q1}), and reports a simulation in which the link between
two groups of firms is made narrower step by step.

The section concerns $\sigma^2_\alpha$, $\sigma^2_\psi$ and the covariance. For the variance of the sum the matrix
$S_{xx}^{-1/2}AS_{xx}^{-1/2}$ equals $(I-uu')/n_{\rm py}$ for a vector $u$ of unit length, so its non-zero
eigenvalues are all equal: the largest accounts for one over the rank of the sum of squared eigenvalues, a
share that vanishes as the number of workers and firms grows. A single direction cannot dominate there.

\subsection{The decomposition of KSS}

Let $\lambda_1$ be the eigenvalue of $S_{xx}^{-1/2}AS_{xx}^{-1/2}$ that is largest in absolute value and $q_1$ its
eigenvector of unit length, and let
\[
w_{i1}=x_i'S_{xx}^{-1/2}q_1,\qquad \hat b_1=\sum_iw_{i1}\,y_i,\qquad \hat\V[\hat b_1]=\sum_iw_{i1}^2\,\hat\sigma_i^2 ,
\]
so that $\sum_iw_{i1}^2=1$. KSS write
$\hat\theta=\lambda_1\big(\hat b_1^{\,2}-\hat\V[\hat b_1]\big)+\hat\theta_{\rm r}$
and base the interval on the joint normality of $\hat b_1$ and the remainder $\hat\theta_{\rm r}$. It needs the three
entries of their covariance matrix $\Sigma$,
\begin{equation}\label{eq:q1cov}
\Sigma_{bb}=\sum_iw_{i1}^2\sigma_i^2,\qquad
\Sigma_{b{\rm r}}=2\sum_iw_{i1}\,m_i^{\rm r}\,\sigma_i^2,\qquad
\Sigma_{\rm rr}=4\sum_i(m_i^{\rm r})^2\sigma_i^2+2\sum_i\sum_{\ell\ne i}(C_{i\ell}^{\rm r})^2\sigma_i^2\sigma_\ell^2 ,
\end{equation}
where $C^{\rm r}$ are the weights of the remainder and $m^{\rm r}=C^{\rm r}\mu$. Let
$\rho=\Sigma_{b{\rm r}}/(\Sigma_{bb}\Sigma_{\rm rr})^{1/2}$ be the correlation of $\hat b_1$ and $\hat\theta_{\rm r}$.

\subsection{The three estimates}

Three facts, all algebra, make the estimates of this note fit:
\begin{enumerate}
\item The remainder is itself a leave-out estimator, $\hat\theta_{\rm r}=\sum_i\sum_{\ell\ne i}C^{\rm r}_{i\ell}y_iy_\ell$,
with weights formed as in \eqref{eq:ustat} from $A^{\rm r}=A-\lambda_1S_{xx}^{1/2}q_1q_1'S_{xx}^{1/2}$ in place of
$A$; that is, from $B^{\rm r}_{i\ell}=B_{i\ell}-\lambda_1w_{i1}w_{\ell1}$.
\item The rows of a stayer have identical regressors, so $w_{i1}$ takes one value $w_{g1}$ on them, and the
weight between two of them is again a constant,
\[
c_g^{\rm r}=\Big(\frac{A_{gg}}{T_g^{\,2}}-\lambda_1w_{g1}^2\Big)\frac{T_g}{T_g-1}.
\]
\item Over the rows of a stayer the leave-out estimates add up to $T_g\,q_g$. So $\hat\V[\hat b_1]$ already uses
the stayers' sample variances.
\end{enumerate}
With $\hat m^{\rm r}=C^{\rm r}y$ and the variance estimates $\st_i^2$ of Section~\ref{sec:estimator} (the stayer's
$q_g$, the fitted model for movers), the estimates are
\begin{align}
\hat\Sigma_{bb}&=\max\Big\{\sum_iw_{i1}^2\,\hat\sigma_i^2,\ \sum_iw_{i1}^2\,\st_i^2\Big\},\qquad
\hat\Sigma_{b{\rm r}}=2\sum_iw_{i1}\,\hat m_i^{\rm r}\,\st_i^2 ,\label{eq:q1est}\\
\hat\Sigma_{\rm rr}&=4\sum_i(\hat m_i^{\rm r})^2\st_i^2-2\sum_i\sum_{\ell\ne i}(C^{\rm r}_{i\ell})^2\st_i^2\st_\ell^2
+4\sum_{g\in\cG}(c_g^{\rm r})^2q_g^2\,\frac{(T_g-2)(T_g-1)}{T_g+1}.\notag
\end{align}
The last is \eqref{eq:Vhat} for the remainder; the second is its analogue for the covariance.

\paragraph{Why the larger of two numbers for $\hat\Sigma_{bb}$.} The first number in the braces is KSS's
$\hat\V[\hat b_1]$. It is unbiased and needs no model for anyone, because $\Sigma_{bb}$ is linear in the error
variances. But with a bottleneck $\hat b_1$ rests on the few rows that link the two groups, and the leave-out
estimates of a few rows are noisy: their weighted sum can come out far too small, or negative, and a
$\hat\Sigma_{bb}$ that is too small makes the interval too short. The second number uses the fitted variances. It
is positive and stable, but only as good as the model on those few rows. Taking the larger guards against the
downward errors of the first without relying on the second. The price is, as a rule, a longer interval
(Corollary~\ref{cor:q1} says when it is certain to be longer). The fitted number is positive whenever the fitted
variances are, which the floors of Section~\ref{sec:estimator} ensure unless residuals vanish; the interval is
reported only if $\hat\Sigma$ is positive definite. The correction inside $\hat\theta_{\rm r}$ is not affected:
the point estimate and the remainder are KSS's.

\subsection{Consistency}

\begin{proposition}[The covariance matrix for $q=1$]\label{prop:q1}
Assume A1 to A3, the moment bounds of Proposition~\ref{prop:conc} and Assumption~1 of KSS, suppose
$\Sigma_{\rm rr}>0$, and give the movers' rows non-random numbers $\st_i^2=\sigma_i^2+\Delta_i$ with
$|\st_i^2|\le\bth\,\sigma_i^2$ for a constant $\bth\ge1$. Let
$\phi_g^{\rm r}$ and $\lambda_{\max}^{\rm r}$ be the quantities of Proposition~\ref{prop:conc} for the remainder,
that is, with $C^{\rm r}$, $m^{\rm r}$ and $\Sigma_{\rm rr}$ in place of $C$, $m$ and $\V[\hat\theta]$, and let
$\epsilon_{\rm r}=\max_g\phi_g^{\rm r}+(\lambda_{\max}^{\rm r})^2/\Sigma_{\rm rr}$. Write
$\hat\Sigma_{bb}^{\rm lo}=\sum_iw_{i1}^2\hat\sigma_i^2$ for the leave-out number in \eqref{eq:q1est},
$d_{\rm r}(\Delta)=2\sum_{i\in\cR}w_{i1}m_i^{\rm r}\Delta_i$, and $b_{\rm r}(\Delta)$ for \eqref{eq:b} formed with
$C^{\rm r}$ and $m^{\rm r}$. Then, in every sample,
\[
\E\hat\Sigma_{bb}^{\rm lo}=\Sigma_{bb},\qquad \E\hat\Sigma_{b{\rm r}}=\Sigma_{b{\rm r}}+d_{\rm r}(\Delta),\qquad
\E\hat\Sigma_{\rm rr}=\Sigma_{\rm rr}+b_{\rm r}(\Delta),
\]
and
\begin{align*}
\E\Big[\Big(\frac{\hat\Sigma_{bb}^{\rm lo}}{\Sigma_{bb}}-1\Big)^2\Big]&\le c_b\,\max_iw_{i1}^2,\\
\E\Big[\frac{\big(\hat\Sigma_{b{\rm r}}-\Sigma_{b{\rm r}}-d_{\rm r}(\Delta)\big)^2}{\Sigma_{bb}\,\Sigma_{\rm rr}}\Big]&\le c_{b{\rm r}}\,\bth^2\,\epsilon_{\rm r},\qquad
\E\Big[\Big(\frac{\hat\Sigma_{\rm rr}-\Sigma_{\rm rr}-b_{\rm r}(\Delta)}{\Sigma_{\rm rr}}\Big)^2\Big]\le c'\,\bth^2\,\epsilon_{\rm r},
\end{align*}
where $c_b$, $c_{b{\rm r}}$ and $c'$ depend only on the constants of the assumptions. Consequently, along a
sequence of designs in which $\max_iw_{i1}^2\to0$, $\epsilon_{\rm r}\to0$, $b_{\rm r}(\Delta)/\Sigma_{\rm rr}\to0$,
$d_{\rm r}(\Delta)/(\Sigma_{bb}\Sigma_{\rm rr})^{1/2}\to0$ and $|\rho|$ stays below a number less than one, the matrix
$\hat\Sigma^{\rm lo}$ with entries $\hat\Sigma_{bb}^{\rm lo}$, $\hat\Sigma_{b{\rm r}}$, $\hat\Sigma_{\rm rr}$ satisfies
$\Sigma^{-1/2}\hat\Sigma^{\rm lo}\Sigma^{-1/2}\to I$ in probability. The same holds when the movers' rows receive
the fitted values of Corollary~\ref{cor:transfer}, with $\Delta_i=v_i^*-\sigma_i^2$ and $n_{\cR}\to\infty$.
\end{proposition}

The proposition is about $\hat\Sigma^{\rm lo}$, the matrix with KSS's first entry. The matrix with the rule
\eqref{eq:q1est} need not be consistent: where the fitted variances overstate $\Sigma_{bb}$ its first entry
stays too large. What the rule keeps is coverage. Let $\hat\rho_{\rm lo}$ be the correlation computed from
$\hat\Sigma^{\rm lo}$.

\begin{corollary}[Coverage]\label{cor:q1}
\textup{(a)} In any sample in which $\hat\Sigma^{\rm lo}$ is positive definite and $\hat\rho_{\rm lo}^2\le2/3$, the
interval of Section~\ref{sec:q1int} formed with \eqref{eq:q1est} contains the interval formed with
$\hat\Sigma^{\rm lo}$.
\textup{(b)} Suppose the conditions of the last part of Proposition~\ref{prop:q1} hold, with $|\rho|$ at most a
constant strictly below $(2/3)^{1/2}$, together with the hypotheses of KSS's Theorem~3 for $q=1$: their
Assumptions~1 and~2, $\max_iw_{i1}^2\to0$, and
$\Sigma_{\rm rr}^{-1}\max_i\big((\tilde x_{i1}'\beta)^2+(\check x_{i1}'\beta)^2\big)\to0$. Then the interval formed
with \eqref{eq:q1est} covers $\theta$ with probability of at least $1-\alpha$ in the limit.
\end{corollary}

Part (a) is algebra (Appendix~\ref{app:proofs}); the bound $2/3$ is on the correlation before the first entry is
raised, and without it a larger first entry can lower the critical value and move an end point inwards. Part (b)
combines three things: KSS's Theorem~3, under which $\hat b_1$ and $\hat\theta_{\rm r}$ are jointly asymptotically
normal; their Lemma~7, under which the interval with a consistent matrix then covers with probability of at
least $1-\alpha$ in the limit; and part (a), which holds with probability tending to one. The condition on
$\tilde x_{i1}$ and $\check x_{i1}$ says that no single row dominates the linear part of the remainder; it is
KSS's and is not implied by the conditions of the proposition. If the fitted variances overstate $\Sigma_{bb}$ in
large samples, the interval is longer than necessary but its coverage bound stands. The conditions of the
proposition, in words:
\begin{itemize}
\item \emph{No single row dominates $\hat b_1$}: $\max_iw_{i1}^2\to0$. This is the Lindeberg condition of KSS's
Theorem~3. KSS read $1/\max_iw_{i1}^2$ as an effective number of observations behind $\hat b_1$.
\item \emph{The remainder has no dominant stayer and no dominant direction}: $\epsilon_{\rm r}\to0$. When the
years per stayer are bounded, a sufficient condition is the remainder's analogue of KSS's conditions, with the
second eigenvalue set against the remaining ones, $\lambda_2^2/\sum_{\ell\ge2}\lambda_\ell^2\to0$
(Appendix~\ref{app:proofs}). This is stronger than KSS's Assumption~2, which sets $\lambda_2^2$ against the sum
over all eigenvalues.
\item \emph{The model's errors average out for the remainder}, in its variance and in its covariance with
$\hat b_1$: $b_{\rm r}(\Delta)/\Sigma_{\rm rr}\to0$ and $d_{\rm r}(\Delta)/(\Sigma_{bb}\Sigma_{\rm rr})^{1/2}\to0$. This is
A4 for the remainder. Both are zero when movers' variances follow the model.
\item \emph{$\hat b_1$ and the remainder are not collinear.} With equal variances $\rho=0$.
\end{itemize}
No independence beyond A1 is used, and the stayers' part is again exact in the mean: only the third condition
involves the model. The zero third moments of A3 are needed for the mean of $\hat\Sigma_{b{\rm r}}$; with
skewed stayers' errors it shifts by a term given in the appendix. As in Section~\ref{sec:results}, the
proposition evaluates the profile at fixed characteristics; the estimator evaluates it at estimated effects.

\subsection{The interval}\label{sec:q1int}

As in KSS (Section~7.1): with $\hat\rho=\hat\Sigma_{b{\rm r}}/(\hat\Sigma_{bb}\hat\Sigma_{\rm rr})^{1/2}$ and the
curvature
\begin{equation}\label{eq:kappa}
\hat\kappa=\frac{2\,|\lambda_1|\,\hat\Sigma_{bb}}{\hat\Sigma_{\rm rr}^{1/2}\,(1-\hat\rho^2)^{1/2}},
\end{equation}
the interval runs from the smallest to the largest value of $\lambda_1b^2+\theta_{\rm r}$ over the ellipse
\[
\begin{pmatrix}b-\hat b_1\\ \theta_{\rm r}-\hat\theta_{\rm r}\end{pmatrix}'\hat\Sigma^{-1}
\begin{pmatrix}b-\hat b_1\\ \theta_{\rm r}-\hat\theta_{\rm r}\end{pmatrix}\le z_{\alpha,\hat\kappa}^{\,2} .
\]
The critical value is that of Andrews and Mikusheva (2016): $z_{\alpha,\kappa}$ is the $1-\alpha$ quantile of
$\big(Z_1^2+(|Z_2|+\kappa^{-1})^2\big)^{1/2}-\kappa^{-1}$ for independent standard normal $Z_1,Z_2$. For
$\alpha=0.05$ it rises from 1.96 at $\kappa=0$ to 2.08 at $\kappa=0.3$, 2.22 at $\kappa=1$, 2.35 at $\kappa=3$ and
2.45 as $\kappa\to\infty$. The curvature measures how much the square of $\hat b_1$ bends the problem. As it
goes to zero the interval approaches the one from a linear approximation,
$\lambda_1\hat b_1^{\,2}+\hat\theta_{\rm r}\pm1.96\,(g'\hat\Sigma g)^{1/2}$ with $g=(2\lambda_1\hat b_1,1)'$. That is
not literally $\hat\theta\pm1.96\,\hat\V[\hat\theta]^{1/2}$: the centre differs by $\lambda_1\hat\V[\hat b_1]$ and the
standard error is built from $\hat\Sigma$. How close the two are in practice is seen below, in the designs
without a bottleneck and on the Veneto panels. The construction needs $\hat\Sigma$ to be positive definite; if
$\hat\Sigma_{bb}$, $\hat\Sigma_{\rm rr}$ or the determinant is not positive, no interval is reported. FEVC's
\texttt{inference(q1)} contains this last step; only the three entries of $\hat\Sigma$ change.

\subsection{Computation}

The eigenvector is found in coefficient space: $v_1=S_{xx}^{-1/2}q_1$ solves $Av_1=\lambda_1S_{xx}v_1$ with
$v_1'S_{xx}v_1=1$. Power iteration on $v\mapsto S_{xx}^{-1}Av$ costs one solve per step, and $w_1=Xv_1$ needs no
further solve. With a bottleneck the leading eigenvalue is well separated and few steps suffice
(\R{q1.steps.bottleneck} in the simulated designs below); without one the iteration can be slow. The decomposition
$\hat\theta=\lambda_1(\hat b_1^{\,2}-\hat\V[\hat b_1])+\hat\theta_{\rm r}$ is an identity for any vector $v_1$ with
$v_1'S_{xx}v_1=1$ and any number $\lambda_1$, so an unconverged iteration does not change the point estimate.
But the results above are for the leading direction, and a small change of $\lambda_1$ from one step to the
next does not certify it; an implementation should check the residual of the eigenvalue equation and say so
when the iteration has not converged. After that $B^{\rm r}_{ii}=B_{ii}-\lambda_1w_{i1}^2$ needs no further
projections, and the remainder is one more component in the calculation of Section~\ref{sec:computation}: one
solve for $\hat m^{\rm r}$ and one for each of the 16 draws: 17 solves per component and 51 for the three, against
the 67 of the standard error, when the draws of the standard error are reused. The steps of the power iteration
come on top. (The prototype does not reuse the draws and spent 85 solves per component.)

The power iteration alone already gives $\lambda_1$, $\max_iw_{i1}^2$ and $\hat\Sigma_{bb}$, and with them a
screening value that needs no further solve,
$\tilde\kappa=2|\lambda_1|\hat\Sigma_{bb}/\hat\V[\hat\theta]^{1/2}$. It is not the curvature and not a bound for it:
its denominator is the whole standard error, which includes the part of $\hat\theta$ that is linear in $\hat b_1$,
so $\tilde\kappa$ is small also when $\hat b_1$ is many standard errors from zero, and its estimate of
$\Sigma_{bb}$ can be too large. In the simulation below its median, over the draws in which
$\hat\V[\hat\theta]$ is positive, was between \R{q1.screen.ratio.min} and \R{q1.screen.ratio.max} times the
curvature. It is a rough guide to whether the interval is worth computing; no
threshold for it has been validated.

\subsection{Simulation}\label{sec:q1sim}

\paragraph{Design.} Synthetic panels with a bottleneck of chosen width: 100 firms in two groups of 50; 12,000
workers observed for two to six years, each attached to one group and assigned a newly drawn firm of that
group with probability 0.10 a year; and $b$ bridging workers, who move once to a firm of the other group. After pruning to the
leave-one-out connected set a panel has about \R{q1.design.rows} rows, of which about \R{q1.design.movers} are
movers' matches. The link is $b=100$, 30 or 10 workers, or there is no bottleneck (one group). Firm effects
differ between the groups by a gap of 0 or 0.3 (their standard deviation within a group is 0.12). Without the
gap $\E\hat b_1$ is \R{q1.b1sd.gap0} standard deviations of $\hat b_1$ from zero, with it \R{q1.b1sd.gap3}. Error
variances are equal (\emph{flat}); a quadratic
in the person effect times a worker-specific factor (\emph{profile}); steeper and three times as high for movers
(\emph{movers}); or three times as high for the bridging workers alone (\emph{bridge}), an indicator that the
model for movers lacks. Errors are normal, or skewed (centred exponential); the skewed designs are outside A3
and are a check of robustness. Appendix~\ref{app:sim} has the details.

The bridging workers carry \R{q1.design.w2bridge} of $\sum_iw_{i1}^2$, and $\max_iw_{i1}^2$ is
\R{q1.maxw2.b100}, \R{q1.maxw2.b30} and \R{q1.maxw2.b10} for $b=100$, 30 and 10: an effective number of
observations behind $\hat b_1$ of about \R{q1.neff.b100}, \R{q1.neff.b30} and \R{q1.neff.b10}. The curvature ranges
from \R{q1.kappa.none.min} without a bottleneck to \R{q1.kappa.b10.max} with ten bridging workers
(\R{q1.kappa.bridge.max} when the bridging workers are noisy).

\paragraph{Procedure.} The estimates and the rule were fixed after a development run on other panels
(\R{q1.dev.draws} draws per design), in which the rule \eqref{eq:q1est} replaced a first version that used the
fitted number only when the leave-out number was unusable. The run reported here uses new panels and new draws,
\R{q1.reps} per design, the estimator of Section~\ref{sec:estimator} throughout (model fitted in every draw,
estimated effects in the profile, 16 draws for the double sums). The designs, the number of draws and the
following tolerances were written down beforehand, for the \R{q1.acc.cells} designs with $b\ge10$ or no
bottleneck: coverage of the 95\% interval of at least 0.92 for each of the three components in each design, and
at least 0.94 on average over these \R{q1.acc.n} cases; no failed draw; no interval missing in more than 1\% of
the draws of a case; coverage of the normal interval for the variance of the sum of at least 0.92 in each
design; and, in each design without a bottleneck and for each component, a mean length of the $q=1$ interval
of at most 1.10 times the mean length of the normal interval. Coverage is a share of all draws: a draw in
which a method gives no interval counts as not covered for that method, and lengths are averaged over the
draws with an interval. The tolerances allow for departures from 0.95; meeting them does not establish nominal
coverage. Two designs with $b=3$ were run as a stress test outside these tolerances: there
$\max_iw_{i1}^2=\R{q1.maxw2.b3}$ and KSS's condition is far from met.

\paragraph{Results.} All \R{q1.draws.total} draws were completed, and every one produced a $q=1$ interval. The
coverage tolerances were met; one tolerance on length was not. Table~\ref{tab:q1} summarises;
Table~\ref{tab:q1cells} in Appendix~\ref{app:sim} lists every design.

\begin{table}[t]\centering\footnotesize\setlength{\tabcolsep}{4.5pt}
\caption{Coverage of 95\% intervals with a bottleneck of $b$ bridging workers (\R{q1.reps} draws per design;
ranges over designs and over $\hat\sigma^2_\alpha$, $\hat\sigma^2_\psi$, $\hat\sigma_{\alpha,\psi}$)}\label{tab:q1}
\tabinput{gen/tab_q1.tex}
\begin{minipage}{0.97\textwidth}\footnotesize\smallskip
\emph{Normal}: $\hat\theta\pm1.96\,\hat\V[\hat\theta]^{1/2}$. \emph{$q=1$}: the interval of this section with
\eqref{eq:q1est}. \emph{Leave-out}: the same with KSS's $\hat\V[\hat b_1]$ alone as $\hat\Sigma_{bb}$; a draw in which
that matrix is not positive definite has no interval and counts as not covered. \emph{True $\Sigma$}: the interval
with the true covariance matrix, which no estimate can use. \emph{Noisy bridge}: the bridging workers' errors
have three times the variance. A normal interval does not exist when $\hat\V[\hat\theta]\le0$; such a draw counts
as not covered. \emph{Length}: mean length of the $q=1$ interval over the mean length of the
normal interval, and over that of the interval with the true matrix, each on the draws in which both exist,
averaged over cases. The Monte Carlo standard error of a coverage of 0.95
is \R{q1.mcse}. The first four blocks have normal errors and the variance patterns flat, profile and movers.
\end{minipage}
\end{table}

\emph{Tolerances.} In the \R{q1.acc.n} cases (three components in \R{q1.acc.cells} designs) the $q=1$ interval
covers between \R{q1.acc.min} and \R{q1.acc.max}, on average \R{q1.acc.mean}; the normal interval for the
variance of the sum covers between \R{q1.sum.min} and \R{q1.sum.max}. The tolerance that failed concerns the
designs without a bottleneck, where the $q=1$ interval was to be at most 10\% longer than the normal one. For
the covariance in one of the three designs it is \R{q1.fail.len} longer; the other eight ratios are between
\R{q1.none.len.others}. In that design the leading direction is a sparse one that rests on a few rows
($\max_iw_{i1}^2=\R{q1.fail.maxw2}$), and the curvature is \R{q1.fail.kappa}, not close to zero: the critical value
alone is \R{q1.fail.zpct} above 1.96, and even with the true matrix the interval is \R{q1.fail.truelen} longer than
the normal one. The tolerance assumed that without a bottleneck the two intervals coincide. That is so only
when the curvature is small. The interval is not wrong there (it covers \R{q1.fail.cover}, the normal one
\R{q1.fail.normalcover}), but the tolerance was set in advance and the result stands as a failure of it.

\emph{The normal interval.} Its coverage falls as the curvature rises: it is \R{q1.bykappa.0.normal} for
$\kappa<0.2$, \R{q1.bykappa.0p2.normal} for $\kappa$ between 0.2 and 0.6, \R{q1.bykappa.0p6.normal} for $\kappa$
between 0.6 and 2, and as low as \R{q1.normal.worst} above that. Three things go wrong. $\hat\theta$ is not close to normal.
The variance estimate is often not positive, so that there is no standard error at all: this happens in
\R{q1.vneg.cases} of the \R{q1.acc.n} cases, in more than a tenth of the draws in \R{q1.vneg.cases10} of them, and in
up to \R{q1.vneg.max} of the draws; such draws count as not covered. (In the case with coverage
\R{q1.normal.worst} it is \R{q1.vneg.worst} of the draws, and among the others the interval covers
\R{q1.normal.worst.cond}; over all cases the coverage among draws with a standard error is at least
\R{q1.normal.cond.min}.) It does not happen for $\kappa<0.6$. And where the standard error exists it can be off,
because with a bottleneck the variance of $\hat\theta$ rests on the bridging rows, where the errors of the model
for movers need not average out: when the bridging workers are noisier than the model can know, the
root-mean-square standard error is \R{q1.se.bridge} times the true one. For the variance of the sum, which has
no dominant direction, the variance estimate is positive in every draw and the standard error is \R{q1.se.sum}
times the true one over all designs.

\emph{The $q=1$ interval.} It covers at least \R{q1.acc.min} in every case, also with skewed errors and with
noisy bridging workers. It is conservative: in \R{q1.acc.above97} of the \R{q1.acc.n} cases it covers more than
0.97. Most of that is the construction of KSS and not the estimates: with the true matrix the coverage is
between \R{q1.true.acc.min} and \R{q1.true.acc.max}. Relative to the interval with the true matrix, the estimated
one is on average \R{q1.lentrue.all.b100}, \R{q1.lentrue.all.b30} and \R{q1.lentrue.all.b10} times as long for
$b=100$, 30 and 10 (the designs of the second to fourth rows of Table~\ref{tab:q1}), and at most
\R{q1.lentrue.acc.max} times in any case.

\emph{The rule for $\hat\Sigma_{bb}$.} With ten bridging workers the leave-out number alone is not usable: it is
not positive, or the matrix is not positive definite, in \R{q1.lonone.b10all} of draws, and the coverage of the
interval built on it is \R{q1.lo.b10all}. Its standard deviation across draws is then \R{q1.bbcv.all.b10} times
the true $\Sigma_{bb}$, against \R{q1.bbcv.all.b100} with a hundred. With thirty or more, or without a bottleneck, the
leave-out number alone covers \R{q1.lo.ge30}, and there the rule costs length without need. The fitted number
alone fails where the model is wrong on the bridging rows: with noisy bridging workers it is
\R{q1.bbfit.bridge} of the true value and the interval built on it covers between \R{q1.fit.bridge.min} and
\R{q1.fit.bridge.max}. The larger of the two does well in both situations.

\emph{The three entries.} Across the \R{q1.acc.cells} designs the mean of the leave-out number is \R{q1.cal.bb.acc} times
$\Sigma_{bb}$, the mean of $\hat\Sigma_{\rm rr}$ is \R{q1.cal.rr.acc} times $\Sigma_{\rm rr}$ (the largest values are
for skewed errors, where A3 fails), and the mean error of $\hat\Sigma_{b{\rm r}}$ is at most \R{q1.cal.br.acc} in
units of $(\Sigma_{bb}\Sigma_{\rm rr})^{1/2}$.

\emph{Stress test.} With three bridging workers the $q=1$ interval covers between \R{q1.q1.b3.min} and
\R{q1.q1.b3.max}, and the leave-out number is unusable in up to \R{q1.lonone.b3.max} of draws. KSS's condition
is not met there, the errors are normal, and nothing should be read into the result beyond these two designs.

\subsection{The Veneto panels}\label{sec:q1real}

The same calculation was run on the six-year Veneto panel of Section~\ref{sec:evidence} and on the two smaller
designs of Section~\ref{sec:validation} (Table~\ref{tab:q1real}). None of them has a dominant direction. The
leading eigenvalue accounts for \R{q1.real.share} of the sum of squared eigenvalues of the three components, the
estimated curvature is \R{q1.real.kappa}, and the $q=1$ interval is \R{q1.real.lenratio} times as long as the
normal one: the two nearly coincide, as they should. For the variance of the sum the share is
\R{q1.real.sum.share}.

The leading direction is a sparse one in these designs too: $\max_iw_{i1}^2$ is \R{q1.real.maxw2}, and half of
$\sum_iw_{i1}^2$ sits on \R{q1.real.half} rows. With so few rows the two numbers in \eqref{eq:q1est} are far apart,
in either direction: the leave-out number is \R{q1.real.bbratio} times the fitted one, and the rule took the
fitted one in \R{q1.real.fitted.n} of the \R{q1.real.n} cases. Because the curvature is small, this hardly moves
the interval. Without a dominant direction the power iteration can be slow: it took \R{q1.real.steps} steps,
against \R{q1.steps.bottleneck} in the simulated bottlenecks, and \R{q1.real.power.seconds} seconds per component
in the prototype. It is stopped at 200 steps, and in \R{q1.real.unconverged} of the \R{q1.real.n} cases it stopped
there without meeting its tolerance; the row of Table~\ref{tab:q1real} with 200 steps is that case, and its
direction is the last iterate, not a certified leading direction.

\begin{table}[t]\centering\footnotesize\setlength{\tabcolsep}{4.5pt}
\caption{The leading direction and the two 95\% intervals on the Veneto designs}\label{tab:q1real}
\tabinput{gen/tab_q1real.tex}
\begin{minipage}{0.97\textwidth}\footnotesize\smallskip
\emph{Steps}: of the power iteration (at most 200). \emph{Ratio}: length of the $q=1$ interval over length of the
normal interval.
\end{minipage}
\end{table}

\subsection{Limits}

The coverage result rests on KSS's, and it requires that no single row dominates $\hat b_1$. Two groups of
firms linked by a handful of movers are not covered by it: adding workers elsewhere in the network does not
make $\max_iw_{i1}^2$ small if the same few bridging workers keep carrying the direction. In the stress test
with three bridging workers the interval still covered in two designs with normal errors; that is not a
guarantee. The simulation shows how the procedure is calibrated on the designs and variance patterns that were
run. It does not establish coverage beyond them, and it does not replace the conditions of
Corollary~\ref{cor:q1}. It uses synthetic panels with one bottleneck; several weak links at once would need
$q>1$, which is not treated here. For $\sigma^2_\psi$ every stayer has $w_{g1}=0$, so stayers do not enter $\hat\Sigma_{bb}$ and
$\hat\Sigma_{b{\rm r}}$, and those two entries rest on the leave-out estimates and the model for movers. For the
covariance and for $\sigma^2_\alpha$ stayers do enter, and can carry much of the remainder's variance; there the
repeat measurements matter. The rule \eqref{eq:q1est} was chosen after a development run, to guard against
intervals that are too short; where the fitted variances overstate the variance of $\hat b_1$ it gives away
length.
